% Problema: discontinuidades evitables (castellano).
%
% \dmarks y \answerspace solo tienen efecto en un examen. En la hoja de
% problemas no ocupan nada, así que el mismo fichero sirve para las dos.

\begin{exercise}[Discontinuidades evitables]
\begin{parts}
  \item Halla el valor de $a$ para el que la función
        \[
          f(x) =
          \begin{cases}
            \dfrac{x^2 + x - 2}{x - 1} & \text{si } x \neq 1, \\[1ex]
            a & \text{si } x = 1,
          \end{cases}
        \]
        es continua en $\mathbb{R}$.
  \item Demuestra que ningún valor de $b$ hace continua en $0$ la función
        $g(x) = 1/x$ para $x \neq 0$, con $g(0) = b$.
\end{parts}

\dmarks{4}
\answerspace{6cm}

\begin{hint}
En (i), factoriza el numerador. En (ii), mira qué le pasa a $|g(x)|$ cuando
$x$ se acerca a $0$.
\end{hint}

\begin{answer}
(i) $a = 3$. (ii) $g$ no tiene límite en $0$: $|1/x|$ se hace tan grande como
se quiera.
\end{answer}

\begin{solution}
\begin{parts}
  \item Para $x \neq 1$, $f$ es un cociente de polinomios cuyo denominador no
        se anula, luego es continua. En $x = 1$, como
        $x^2 + x - 2 = (x - 1)(x + 2)$,
        \[
          \lim_{x \to 1} f(x) = \lim_{x \to 1} (x + 2) = 3 ,
        \]
        así que $f$ es continua en $1$ si y solo si $a = 3$.
  \item Si $g$ fuera continua en $0$, con $\varepsilon = 1$ habría un
        $\delta > 0$ tal que $|1/x - b| < 1$, y por tanto
        $|1/x| < |b| + 1$, para todo $0 < |x| < \delta$. Pero
        $x = \min\{\delta/2, 1/(|b| + 1)\}$ cumple $0 < x < \delta$ y
        $1/x \geq |b| + 1$. Contradicción.
\end{parts}
\end{solution}

\begin{marking}
(i) Dos puntos: uno por el límite y otro por concluir $a = 3$ justificando la
continuidad fuera de $1$. (ii) Dos puntos. Decir solo «el límite es infinito»
sin justificarlo vale medio punto.
\end{marking}
\end{exercise}
